\subsection{শেষ অনুশীলন: নতুন টেবিলে নিজে কুয়েরি লেখো}


\subsubsection{অনুশীলনের টেবিল--১: Worker}


\begin{flushleft}

\small

\begin{tabular}{|l|l|l|l|l|l|}
\hline

\chSixHeader
WorkerID & Name & Department & Salary & City & JoinYear\\
\hline
W01 & Rakib & Production & 28000 & Dhaka & 2022\\\hline
W02 & Sumaiya & Accounts & 35000 & Rajshahi & 2021\\\hline
W03 & Imran & Production & 32000 & Khulna & 2023\\\hline
W04 & Tania & Sales & 42000 & Dhaka & 2020\\\hline
W05 & Sakib & Sales & 30000 & Rajshahi & 2022\\\hline
W06 & Mahi & Accounts & 38000 & Dhaka & 2021\\
\hline

\end{tabular}

\end{flushleft}


\begin{chSixSideBySide}[.48]
\begin{enumerate}

\item Worker টেবিলের সব কর্মীর তথ্য দেখাও।

\item শুধু Name ও Salary দেখাও।

\item যেসব কর্মীর City-এর মান Dhaka, তাদের তথ্য দেখাও।

\item যেসব কর্মীর Salary 30000-এর বেশি, তাদের তথ্য দেখাও।

\item যেসব কর্মীর Department-এর মান Production অথবা Sales, তাদের তথ্য দেখাও।

\end{enumerate}
\chSixNextColumn
\begin{enumerate}[start=6]
\item যেসব কর্মীর Salary 28000 থেকে 35000-এর মধ্যে, তাদের তথ্য দেখাও।

\item যেসব কর্মীর JoinYear-এর মান 2022, তাদের তথ্য দেখাও।

\item Salary বড় থেকে ছোট ক্রমে সাজাও।

\item Department-এর মান Accounts এমন কর্মীদের গড় Salary বের করো।

\item সর্বোচ্চ Salary বের করো।

\end{enumerate}
\end{chSixSideBySide}


\subsubsection{অনুশীলনের টেবিল--২: Product}


\begin{flushleft}

\small

\begin{tabular}{|l|l|l|l|l|}
\hline

\chSixHeader
ProductID & ProductName & Category & Price & Stock\\
\hline
P01 & Keyboard & Electronics & 1200 & 15\\\hline
P02 & Mouse & Electronics & 650 & 8\\\hline
P03 & Notebook & Stationery & 120 & 40\\\hline
P04 & Pen Set & Stationery & 250 & 5\\\hline
P05 & Headphone & Electronics & 1800 & 12\\
\hline

\end{tabular}

\end{flushleft}


\begin{multicols}{2}
\begin{enumerate}

\item যেসব পণ্যের Stock 10-এর কম, তাদের তথ্য দেখাও।

\item যেসব পণ্যের Price 500 থেকে 1500-এর মধ্যে, তাদের তথ্য দেখাও।

\item যেসব পণ্যের Category-এর মান Electronics, তাদের তথ্য দেখাও।

\item মোট Stock বের করো।

\item Lowest Price বের করো।

\item Highest Price বের করো।

\item Price Descending Order-এ সাজাও।

\item যেসব পণ্যের ProductName H দিয়ে শুরু, তাদের তথ্য দেখাও।

\end{enumerate}
\end{multicols}


\subsubsection{অনুশীলনের টেবিল--৩: Course}


\begin{flushleft}

\small

\begin{tabular}{|l|l|l|l|}
\hline

\chSixHeader
CourseID & CourseName & Credit & Teacher\\
\hline
C101 & Database & 3 & Salma\\\hline
C102 & Programming & 4 & Rahat\\\hline
C103 & Networking & 3 & Farid\\\hline
C104 & Web Design & 3 & Salma\\\hline
C105 & AI Basics & 4 & Nabila\\
\hline

\end{tabular}

\end{flushleft}


\begin{multicols}{2}
\begin{enumerate}

\item যেসব কোর্সের Credit-এর মান 4, তাদের তথ্য দেখাও।

\item Salma যে কোর্সগুলো পড়ান, সেগুলোর তথ্য দেখাও।

\item যেসব কোর্সের CourseName D দিয়ে শুরু, তাদের তথ্য দেখাও।

\item মোট Course সংখ্যা বের করো।

\item Teacher অনুযায়ী Course সংখ্যা বের করো।

\item CourseName Ascending Order-এ সাজাও।

\end{enumerate}
\end{multicols}


% ============================================================
% PRACTICE HOSPITAL
% ============================================================


\subsubsection{অনুশীলনের টেবিল--৪: Hospital}


\begin{flushleft}

\small

\begin{tabular}{|l|l|l|l|l|}
\hline

\chSixHeader
PatientID & Name & Age & BloodGroup & Bill\\
\hline
H01 & Rafi & 35 & A+ & 12000\\\hline
H02 & Nila & 22 & O+ & 8000\\\hline
H03 & Hasan & 60 & AB+ & 25000\\\hline
H04 & Mitu & 48 & B+ & 18000\\\hline
H05 & Sami & 65 & O+ & 30000\\
\hline

\end{tabular}

\end{flushleft}


\begin{multicols}{2}
\begin{enumerate}

\item যেসব রোগীর Age 50-এর বেশি, তাদের তথ্য দেখাও।

\item যেসব রোগীর রক্তের গ্রুপ O+, তাদের তথ্য দেখাও।

\item যেসব রোগীর Bill 10000 থেকে 25000-এর মধ্যে, তাদের তথ্য দেখাও।

\item Highest Bill বের করো।

\item Average Bill বের করো।

\item Bill Descending Order-এ সাজাও।

\item যেসব রোগীর Name S দিয়ে শুরু, তাদের তথ্য দেখাও।

\end{enumerate}
\end{multicols}


% ============================================================
% FINAL PROBLEM SOLVING MAP
% ============================================================



\subsection{কুয়েরির ধারণা যাচাই}
\begin{practicebox}
\textbf{নিজে যাচাই করো।} কুয়েরি ও SQL অংশের মূল Student ও FeeList সারণি ব্যবহার করো। (১) Student সারণিতে COUNT(*) ও COUNT(Email) কেন আলাদা? (২) FeeList-এ WHERE TuitionFee>1000 এবং HAVING AVG(TuitionFee)>1100 একই ফল বাছাই করে কি?

\end{practicebox}

\subsection{আরও দুটি বহুনির্বাচনি প্রয়োগ}
\begin{enumerate}
\item Science তিনবার, Humanities দুবার এবং Business Studies তিনবার থাকলে \texttt{SELECT DISTINCT Department FROM Student;} কয়টি সারি দেখাবে?\par
ক. 2\quad খ. 3\quad গ. 5\quad ঘ. 8\par
\chSixMCQAnswer{খ}
\item CREATE, ALTER ও DROP-কে কোন শ্রেণিতে রাখা হয়?\par
ক. DDL\quad খ. DCL\quad গ. TCL\quad ঘ. QBE\par
\chSixMCQAnswer{ক}
\end{enumerate}


\subsection{QBE-তে শর্তের অবস্থান: প্রয়োগমূলক MCQ}
\begin{enumerate}
\item \begin{minipage}[t]{\linewidth}Department-এর Criteria সারিতে Science, কিন্তু GPA-এর Or সারিতে \texttt{>=4.80} দেওয়া আছে। কোন রেকর্ড নির্বাচিত হবে?
\begin{tabularx}{\linewidth}{@{}Y Y@{}}
ক. শুধু Science ও GPA অন্তত 4.80 & খ. Science অথবা GPA অন্তত 4.80\\
গ. শুধু Science-এর বাইরের & ঘ. শুধু GPA ঠিক 4.80\\
\end{tabularx}
\chSixMCQAnswer{খ}
\end{minipage}
\item \begin{minipage}[t]{\linewidth}Department-এর Show বন্ধ, কিন্তু Criteria-তে Science আছে। কোনটি সঠিক?
\begin{tabularx}{\linewidth}{@{}Y Y@{}}
ক. শর্ত কাজ করবে না & খ. মূল কলাম মুছে যাবে\\
গ. শর্ত কাজ করবে, কলাম দেখাবে না & ঘ. সব বিভাগ আসবে\\
\end{tabularx}
\chSixMCQAnswer{গ}
\end{minipage}
\end{enumerate}
